% Options for packages loaded elsewhere
\PassOptionsToPackage{unicode}{hyperref}
\PassOptionsToPackage{hyphens}{url}
\documentclass[
  10pt,
  a4paper,
]{article}
\usepackage{xcolor}
\usepackage[margin=22mm]{geometry}
\usepackage{amsmath,amssymb}
\setcounter{secnumdepth}{-\maxdimen} % remove section numbering
\usepackage{iftex}
\ifPDFTeX
  \usepackage[T1]{fontenc}
  \usepackage[utf8]{inputenc}
  \usepackage{textcomp} % provide euro and other symbols
\else % if luatex or xetex
  \usepackage{unicode-math} % this also loads fontspec
  \defaultfontfeatures{Scale=MatchLowercase}
  \defaultfontfeatures[\rmfamily]{Ligatures=TeX,Scale=1}
\fi
\usepackage{lmodern}
\ifPDFTeX\else
  % xetex/luatex font selection
\fi
% Use upquote if available, for straight quotes in verbatim environments
\IfFileExists{upquote.sty}{\usepackage{upquote}}{}
\IfFileExists{microtype.sty}{% use microtype if available
  \usepackage[]{microtype}
  \UseMicrotypeSet[protrusion]{basicmath} % disable protrusion for tt fonts
}{}
\makeatletter
\@ifundefined{KOMAClassName}{% if non-KOMA class
  \IfFileExists{parskip.sty}{%
    \usepackage{parskip}
  }{% else
    \setlength{\parindent}{0pt}
    \setlength{\parskip}{6pt plus 2pt minus 1pt}}
}{% if KOMA class
  \KOMAoptions{parskip=half}}
\makeatother
\setlength{\emergencystretch}{3em} % prevent overfull lines
\providecommand{\tightlist}{%
  \setlength{\itemsep}{0pt}\setlength{\parskip}{0pt}}
\usepackage{mathrsfs}
\usepackage{mathtools}
\usepackage{xurl}
\emergencystretch=3em
\setcounter{tocdepth}{1}
\newcommand{\Beta}{\mathrm{B}}
\DeclareUnicodeCharacter{00B2}{\ensuremath{^2}}
\DeclareUnicodeCharacter{00B1}{\ensuremath{\pm}}
\DeclareUnicodeCharacter{00D7}{\ensuremath{\times}}
\DeclareUnicodeCharacter{2192}{\ensuremath{\to}}
\DeclareUnicodeCharacter{21D2}{\ensuremath{\Rightarrow}}
\DeclareUnicodeCharacter{21A6}{\ensuremath{\mapsto}}
\DeclareUnicodeCharacter{2208}{\ensuremath{\in}}
\DeclareUnicodeCharacter{2209}{\ensuremath{\notin}}
\DeclareUnicodeCharacter{2212}{\ensuremath{-}}
\DeclareUnicodeCharacter{2227}{\ensuremath{\wedge}}
\DeclareUnicodeCharacter{2228}{\ensuremath{\vee}}
\DeclareUnicodeCharacter{2260}{\ensuremath{\ne}}
\DeclareUnicodeCharacter{2264}{\ensuremath{\le}}
\DeclareUnicodeCharacter{2265}{\ensuremath{\ge}}
\DeclareUnicodeCharacter{2297}{\ensuremath{\otimes}}
\DeclareUnicodeCharacter{00B3}{\ensuremath{^3}}
\DeclareUnicodeCharacter{2074}{\ensuremath{^4}}
\DeclareUnicodeCharacter{2076}{\ensuremath{^6}}
\DeclareUnicodeCharacter{03BB}{\ensuremath{\lambda}}
\DeclareUnicodeCharacter{03BC}{\ensuremath{\mu}}
\DeclareUnicodeCharacter{03B5}{\ensuremath{\epsilon}}
\DeclareUnicodeCharacter{03C4}{\ensuremath{\tau}}
\DeclareUnicodeCharacter{03A6}{\ensuremath{\Phi}}
\DeclareUnicodeCharacter{03B1}{\ensuremath{\alpha}}
\DeclareUnicodeCharacter{03C3}{\ensuremath{\sigma}}
\DeclareUnicodeCharacter{03B4}{\ensuremath{\delta}}
\DeclareUnicodeCharacter{03C0}{\ensuremath{\pi}}
\DeclareUnicodeCharacter{039B}{\ensuremath{\Lambda}}
\DeclareUnicodeCharacter{2202}{\ensuremath{\partial}}
\DeclareUnicodeCharacter{0393}{\ensuremath{\Gamma}}
\DeclareUnicodeCharacter{2205}{\ensuremath{\varnothing}}
\DeclareUnicodeCharacter{221E}{\ensuremath{\infty}}
\DeclareUnicodeCharacter{2200}{\ensuremath{\forall}}
\DeclareUnicodeCharacter{00B9}{\ensuremath{^1}}
\DeclareUnicodeCharacter{00B0}{\ensuremath{{}^\circ}}
\DeclareUnicodeCharacter{211D}{\ensuremath{\mathbb{R}}}
\usepackage{bookmark}
\IfFileExists{xurl.sty}{\usepackage{xurl}}{} % add URL line breaks if available
\urlstyle{same}
\hypersetup{
  hidelinks,
  pdfcreator={LaTeX via pandoc}}

\author{}
\date{}

\begin{document}

{
\setcounter{tocdepth}{3}
\tableofcontents
}
\section{Reconstruction over the full support
lattice}\label{reconstruction-over-the-full-support-lattice}

This calculation starts from the programme's actual support-amplitude
semiring. It proves the precise comparison with the earlier scalar
calculations, identifies information lost by all separate Boolean branch
maps after linearization, and constructs a support-valued extension of
the finite Weil packet form. It does not assert that a zeta zero off the
critical line exists, or that positivity for the classical Weil
distribution has been proved.

The source definition and the existing spectrum theorem are in
\emph{Split Support Geometry: Universal Zero Fibres, Arithmetic Curves,
and Frobenius-Perfect Quantization}, The Clankers, June 2026, Version
11. The file read for this derivation is
\texttt{split\_\allowbreak{}support\_\allowbreak{}geometry\_\allowbreak{}arithmetic\_\allowbreak{}curve\_\allowbreak{}v11910.\allowbreak{}tex}, SHA256
\texttt{8443cc04\allowbreak{}01b18d37\allowbreak{}3939d992\allowbreak{}f0a2f538\allowbreak{}4ee2fa11\allowbreak{}30f05975\allowbreak{}3d5f3972\allowbreak{}eec5052d},
Definition \texttt{def:\allowbreak{}lattice-\allowbreak{}split}, Theorems
\texttt{thm:\allowbreak{}lattice-\allowbreak{}normal-\allowbreak{}form},
\texttt{thm:\allowbreak{}birkhoff-\allowbreak{}support-\allowbreak{}coordinates},
\texttt{thm:\allowbreak{}lattice-\allowbreak{}prime-\allowbreak{}classification},
\texttt{thm:\allowbreak{}spectral-\allowbreak{}ordinal-\allowbreak{}sum}, and \texttt{thm:\allowbreak{}support-\allowbreak{}base-\allowbreak{}change},
source lines1180--1871. Those results are prior programme mathematics.
The different file
\texttt{split\_\allowbreak{}support\_\allowbreak{}geometry\_\allowbreak{}arithmetic\_\allowbreak{}curve\_\allowbreak{}v11.\allowbreak{}tex}, SHA256
\texttt{5fa6a55a\allowbreak{}a9d1a2e5\allowbreak{}7599e206\allowbreak{}087d3d24\allowbreak{}66b25dd9\allowbreak{}9fa1a90d\allowbreak{}5f86367a\allowbreak{}afab1b4d},
has the explicit branch-spectrum theorem at lines1925--1985. The
editions are not treated as byte-identical.

\subsection{1. The base and its split comparison
maps}\label{the-base-and-its-split-comparison-maps}

Let R be a nonzero commutative unital ring and L a nontrivial bounded
distributive lattice. Write its bounds as 0\_L and 1\_L. Set \[
S_L=G_L(R)=\{z_\lambda=(0,\lambda):\lambda\in L\}
\cup\{\widehat r=(r,1_L):r\in R\},
\quad z_{1_L}=\widehat0=e_L,\quad z_{0_L}=\tau_L.
\tag{FSR1}
\] The overlap in this union consists of the one element e\_L. The
operations are \[
(r,\lambda)+(s,\mu)=(r+s,\lambda\vee\mu),\qquad
(r,\lambda)(s,\mu)=(rs,\lambda\wedge\mu).
\tag{FSR2}
\] In particular, \(\widehat r+z_\lambda=\widehat r\) and
\(\widehat r z_\lambda=z_\lambda\), including r=0. Closure follows
because every nonzero amplitude has top support. The semiring laws
follow from those of R and the distributive lattice law. Its zero is
tau\_L and its unit is widehat1.

Put S=G\_B(R), with B=\{0,1\}. The endpoint inclusion and each Boolean
branch give maps \[
\iota:S\longrightarrow S_L,\quad \tau\mapsto\tau_L,\quad r\mapsto\widehat r,
\qquad
\beta_\epsilon:S_L\longrightarrow S,\quad
z_\lambda\mapsto z_{\epsilon(\lambda)},\quad \widehat r\mapsto r,
\qquad \beta_\epsilon\iota=\mathrm{id}_S.
\tag{FSR3}
\] Here epsilon:L→B preserves the bounds, joins, and meets. Both maps
preserve every operation in FSR2: on the zero fibre this is the
lattice-map identity; on two supported amplitudes it is the identity of
R; on a mixed pair it follows from the two identities after FSR2. Thus
iota is injective and every beta\_epsilon is surjective. The kernel
congruence of beta\_epsilon identifies exactly the zero-fibre elements
with the same epsilon value, and no two distinct nonzero amplitudes. Its
inverse image of the actual scalar zero is \[
\beta_\epsilon^{-1}(\tau)=\{z_\lambda:\epsilon(\lambda)=0\}.
\tag{FSR4}
\] Its inverse image of the supported scalar zero is the complementary
set of zero-fibre elements. This is the exact distinction between a zero
ideal and a congruence in this comparison.

The amplitude projection p\_L:S\_L→R has zero fibre Z\_L=\{z\_lambda\}.
Every homomorphism to a ring kills Z\_L: the equality
e\_L+z\_lambda=e\_L implies this by additive cancellation in the target.
The restriction to the supported copy then gives a unique ring
homomorphism R→A. This proves the universal property of p\_L, without
identifying e\_L with tau\_L in S\_L.

For a domain R, the supported zero is a prime element in the
divisibility sense: \[
e_L S_L=Z_L,
\qquad e_L\mid xy\ \Longleftrightarrow\ p_L(x)p_L(y)=0
\ \Longleftrightarrow\ e_L\mid x\ \text{or}\ e_L\mid y.
\tag{FSR5}
\] The first equality follows from e\_L z\_lambda=z\_lambda and e\_L
widehat r=e\_L. The next equivalence follows from the first equality and
the multiplication law. The last uses precisely the absence of zero
divisors in R. The ideal Z\_L is proper because widehat1 is outside it.
The element e\_L is not a unit, since every product with it has
amplitude zero. It is reducible, since e\_L=e\_L e\_L with two nonunits.
Its primality and its multiplicative idempotence are simultaneous facts.

For completeness the entire semiring-prime comparison can be obtained
directly. If an ideal contains a supported amplitude, multiplication by
e\_L puts e\_L in it, hence puts all of Z\_L in it. Its supported
amplitudes form a ring ideal J. An ideal containing no supported
amplitude is Z\_a for a proper lattice ideal a: joins follow from
addition and downward closure follows from z\_mu z\_lambda=z\_mu when
mu≤lambda. Thus all ideals are Z\_a and \[
Q_J=Z_L\cup\{\widehat r:r\in J\}.
\tag{FSR6}
\] Testing FSR2 shows that Z\_a is prime exactly when a is a prime
lattice ideal, and Q\_J is prime exactly when J is a ring prime. Every
support prime is strictly below every arithmetic prime. For R=Z this
includes \[
Z_{\mathfrak a}\subsetneq Z_L=(e_L)=Q_{(0)}\subsetneq Q_{(p)}.
\tag{FSR7}
\] The basic opens are D(z\_lambda)=\{Z\_a:lambda∉a\} and D(widehat
r)=Spec\_lat(L) union D\_R(r). These formulas prove the topological
ordinal sum, including its topology. Contraction along beta\_epsilon
sends the scalar prime \{tau\} to Z\_(epsilon inverse0) and Q\_q to
Q\_q. Contraction along iota sends every Z\_a to \{tau\} and Q\_q to
Q\_q. Consequently the scalar spectrum is a retract with an explicitly
chosen section; it is not the entire support spectrum.

\subsection{2. Joint reconstruction before
linearization}\label{joint-reconstruction-before-linearization}

Now assume L finite and let J(L) be its nonzero join-irreducible
elements. For j∈J(L) put theta\_j(lambda)=1 if j≤lambda and 0 otherwise.
A join-irreducible is join-prime: if j≤lambda∨mu, distributivity gives
j=(j∧lambda)∨(j∧mu), so one term equals j. Therefore theta\_j is a
bounded lattice homomorphism. Every nonzero element is the join of the
join-irreducibles below it, by induction in the finite partially ordered
set. Hence these characters jointly distinguish elements of L.

Conversely, a Boolean branch has a prime filter whose least element j is
the meet of its elements. If j=x∨y with neither x nor y equal to j, the
primality of the filter contradicts the minimality of j. Thus j is
join-irreducible, and the branch is theta\_j. This proves the entire
branch list.

The product of all branch maps is the injective semiring map \[
\Beta:S_L\longrightarrow\prod_{j\in J(L)}S,\qquad
\Beta(x)=(\beta_j(x))_j.
\tag{FSR8}
\] Its image is exactly the following set. All coordinate amplitudes
have a common value r. If r≠0, all coordinates are the same supported
element r. If r=0, each coordinate is tau or e, and the set of
coordinates equal to e is an order ideal in J(L). Indeed it is
\{j:j≤lambda\} for the source z\_lambda. Conversely, an order ideal A
gives lambda=join A; join-primality shows that \{j:j≤lambda\}=A. This
proves both surjectivity onto the displayed image and injectivity.

For L=B\^{}d the poset J(L) is an antichain of d atoms. Every subset is
an order ideal, so FSR8 becomes an isomorphism \[
G_{B^d}(R)\ \cong\ \underbrace{G(R)\times_R\cdots\times_R G(R)}_{d\text{ factors}},
\tag{FSR9}
\] where the maps to R are amplitude projections. This is different from
the product of independent amplitude coordinates: the common-amplitude
condition is part of the isomorphism, not an omitted restriction.

For a general finite distributive L, FSR8 is the subsemiring of that
fibre product whose zero-amplitude masks satisfy the stated order-ideal
relations. The additive and multiplicative laws remain FSR2. All earlier
scalar identities pull back along each beta\_j. Conversely, if two
S\_L-valued expressions have equal images under every beta\_j, FSR8
proves equality in S\_L. This last conclusion concerns semiring-valued
expressions; it does not yet concern their free additive envelopes.

\subsection{3. The exact extra kernel created by
linearization}\label{the-exact-extra-kernel-created-by-linearization}

For a commutative coefficient ring k define the contracted meet algebra
\[
C_L=k[(L,\wedge)]/([0_L]).
\tag{FSR10}
\] It is free as a k-module on u\_lambda={[}lambda{]} for lambda≠0\_L,
with u\_0=0, unit u\_1, and u\_lambda u\_mu=u\_(lambda∧mu). It retains
multiplication from support and introduces its own additive linear
combinations. In particular u\_(lambda∨mu) is not defined to be
u\_lambda+u\_mu.

For every nonzero a∈L, whether join-irreducible or not, define \[
\eta_a(u_\lambda)=\mathbf1_{\{a\leq\lambda\}}.
\tag{FSR11}
\] This is a unital k-algebra homomorphism: a≤lambda∧mu holds exactly
when both inequalities hold. With any linear extension of the order, its
matrix has entries 1\_(a≤lambda), is triangular with diagonal1, and is
invertible over the integers. Consequently \[
\eta:C_L\xrightarrow{\sim}\prod_{a\in L\setminus\{0\}}k.
\tag{FSR12}
\] No division in k is used.

An explicit inverse uses the Möbius function of the finite order. Define
mu(a,a)=1 and, for a\textless b, mu(a,b)=-sum\_(a≤c\textless b)mu(a,c).
The upper triangular incidence matrix Z\_(ab)=1\_(a≤b) has inverse with
these entries: the recurrence proves MZ=I, and finite triangular
matrices then also satisfy ZM=I. Define \[
E_a=\sum_{\lambda\leq a}\mu(\lambda,a)u_\lambda.
\tag{FSR13}
\] Then eta\_b(E\_a)=sum\_(b≤lambda≤a)mu(lambda,a)=delta\_(ab). The term
lambda=0 contributes zero. Thus \[
E_aE_b=\delta_{ab}E_a,\qquad
\sum_{a\neq0}E_a=u_1,\qquad
u_\lambda=\sum_{0<a\leq\lambda}E_a.
\tag{FSR14}
\] These identities follow by applying the isomorphism FSR12, so hold
over every coefficient ring.

The linearized Boolean branch theta\_j is precisely eta\_j. Write \[
K_L=\bigcap_{j\in J(L)}\ker\eta_j
=\bigoplus_{a\in L\setminus(\{0\}\cup J(L))}kE_a.
\tag{FSR15}
\] This proves the full kernel, not only its dimension. Over a field its
dimension is \textbar L\textbar−1−\textbar J(L)\textbar. It is a direct
product of copies of k, with multiplicatively idempotent coordinate
projectors. It is zero exactly when every nonzero lattice element is
join-irreducible, equivalently when L is a chain. To verify the last
equivalence, a chain has that property. If L is not a chain, choose
incomparable a,b; then a∨b is nonzero and join-reducible.

For L=B\^{}2=\{0,a,b,1\}, with a∧b=0 and a∨b=1, \[
E_a=u_a,\qquad E_b=u_b,\qquad
w=E_1=u_1-u_a-u_b,
\quad w^2=w\neq0,
\quad\eta_a(w)=\eta_b(w)=0.
\tag{FSR16}
\] Indeed the three terms are distinct basis vectors, proving nonzero,
and direct expansion gives
\((u_1-u_a-u_b)^2=u_1+u_a+u_b-2u_a-2u_b=u_1-u_a-u_b\). Thus the semiring
branch maps in FSR8 distinguish every support element, but their
separately linearized maps omit w. This is a proved failure of
linearized reconstruction from separate branches. Its retained object is
K\_L with inclusion, coordinate projectors, and quotient given in FSR15.

The nonzero a that are not join-irreducible give genuine mixed
observables. For a with join-irreducible decomposition a=join\_(j≤a)j,
\[
\eta_a(u_\lambda)=\prod_{j\leq a}\theta_j(\lambda).
\tag{FSR17}
\] Both sides are1 exactly when every j≤a lies below lambda,
equivalently when a≤lambda. Thus a joint product of branch tests
restores this missing coordinate. Separate branch expectations do not
determine expectations of these products. FSR12 gives the complete
recovery map using all these joint tests.

\subsection{4. Contracted multiplicative coefficients and the scalar
retract}\label{contracted-multiplicative-coefficients-and-the-scalar-retract}

Let Gamma\_L=k{[}(S\_L, multiplication){]}/({[}tau\_L{]}). Put
E={[}e\_L{]} and F=1−E. Multiplication by E sends {[}z\_lambda{]} to
itself and {[}widehat r{]} to {[}e\_L{]}. Thus E Gamma\_L is C\_L,
embedded by u\_lambda↦{[}z\_lambda{]}. The complementary factor F
Gamma\_L has basis \[
v_r=[\widehat r]-[e_L],\qquad r\in R\setminus\{0\},
\tag{FSR18}
\] and v\_0=0. Their multiplication is v\_r v\_s=v\_(rs), with zero when
rs=0. Indeed expansion gives
{[}widehat(rs){]}−{[}e\_L{]}−{[}e\_L{]}+{[}e\_L{]}. The two sets
\{{[}z\_lambda{]}:lambda≠0\} and \{v\_r:r≠0\} form a basis: each
original nonzero supported basis element is v\_r+{[}e\_L{]}, and
comparison of the distinct {[}widehat r{]} coefficients proves linear
independence. Hence \[
\Gamma_L\cong C_L\times D_R,
\qquad D_R=k[(R,\times)]/([0_R]).
\tag{FSR19}
\] The formulas prove this for rings with zero divisors as well as
domains. For R=Z the second factor is \[
D_{\mathbb Z}\cong k[t,X_p:p\text{ rational prime}]/(t^2-1),
\quad v_{-1}\mapsto t,\quad v_p\mapsto X_p.
\tag{FSR20}
\] Unique prime factorization of every nonzero integer proves a
bijection of the displayed monomial bases and compatibility with
multiplication.

For scalar support the same construction is Gamma\_B=k×D\_R. Under FSR19
each branch and the scalar inclusion become \[
\Gamma(\beta_j):(c,d)\longmapsto(\eta_j(c),d),
\qquad
\Gamma(\iota):(a,d)\longmapsto(a\,u_1,d).
\tag{FSR21}
\] These formulas follow on z\_lambda and widehat r, which generate the
algebras. The composite is the identity. The intersection of all branch
kernels is exactly K\_L×0. Thus the earlier arithmetic prime-label
algebra D\_R is preserved as an actual direct factor, and the missing
support coefficient object is explicitly identified. This is stronger
than saying that earlier work might survive.

The evaluation map Gamma\_L→R when k=Z sends {[}widehat r{]} to r and
every {[}z\_lambda{]} to0. For R=Z, FSR19--20 identify its kernel as \[
C_L\times (t+1,\ X_p-p\text{ for every prime }p)\subset C_L\times D_{\mathbb Z}.
\tag{FSR22}
\] The quotient of D\_Z by the displayed ideal is Z: every polynomial
reduces to its integer evaluation, and that evaluation is the inverse of
the resulting inclusion of Z. This proves equality of the kernel ideal.
In particular all of C\_L, including K\_L, is killed by arithmetic
addition. It remains present before that quotient. The inclusion and
quotient in FSR19 and FSR22 give its exact relation to the arithmetic
factor.

\subsection{5. Extension of complexes, operators, and all primary
jets}\label{extension-of-complexes-operators-and-all-primary-jets}

Take k=C for this section. Let (V\^{}bullet,d) be any specified complex
of complex vector spaces, with its original differentials. Its full
support extension is \[
V_L^\bullet=C_L\otimes_{\mathbb C}V^\bullet
=\bigoplus_{a\neq0}E_a\otimes V^\bullet,
\qquad d_L(E_a\otimes v)=E_a\otimes dv.
\tag{FSR23}
\] This is a defined construction, not an identification with a
prismatic complex. Decomposition FSR14 proves directly \[
\ker d_L=\bigoplus_a E_a\otimes\ker d,
\qquad\operatorname{im}d_L=\bigoplus_a E_a\otimes\operatorname{im}d,
\qquad H^q(V_L)=\bigoplus_a E_a\otimes H^q(V).
\tag{FSR24}
\] The first two equalities follow by equating the coefficients of the
independent E\_a. The third is their quotient. Thus a nonzero cohomology
class of a specified old complex remains nonzero in every individual
support sector. For every linear chain map T, the map 1⊗T preserves all
sectors. A homotopy dH+Hd=T−T' extends to 1⊗H and satisfies the
identical formula. The maps eta\_j⊗id are chain maps, with joint kernel
K\_L⊗V\^{}bullet. The scalar inclusion v↦u\_1⊗v is split by every
eta\_j⊗id. These formulas specify exactly which part of this statement
is functorial and which prismatic comparison still requires identifying
the actual input complex.

For the finite packet take the four distinct points \[
\alpha_1=\rho,\qquad\alpha_2=1-\overline\rho,\qquad
\alpha_3=\overline\rho,\qquad\alpha_4=1-\rho,
\] where \(\rho=1/2+\delta+i\gamma\), \(0<\delta<1/2\) and \(\gamma>2\).
Put \(d(s)=\prod_{i=1}^4(s-\alpha_i)\), \(h=d^m\) with \(m\ge1\), and
\(A_h=\mathbb C[s]/(h)\). These parameters define a finite algebra for
every such \(\rho\); they do not assert that \(\rho\) is a zeta zero.
Pairwise coprimality gives the isomorphism to the four rings
\(\mathbb C[t_i]/(t_i^m)\) by Taylor coefficients. To see its inverse
explicitly, take \(h_i=h/(s-\alpha_i)^m\) and multiply \(h_i\) by the
degree-\((m-1)\) Taylor polynomial of \(1/h_i\) at \(\alpha_i\). The
resulting class \(e_i\) is 1 modulo \((s-\alpha_i)^m\) and 0 modulo each
other primary power. Then \[
\sum_{i=1}^4\sum_{b=0}^{m-1}c_{i,b}e_i(s-\alpha_i)^b
\tag{FSR25}
\] is the inverse image of the four Taylor polynomials. Injectivity
follows because a polynomial vanishing to order m at all four points is
divisible by h. This proves the complete primary decomposition and
retains every jet.

Evaluation \(\operatorname{ev}:A_h\to\mathbb C^4\) is surjective with
kernel \(J_h=(d)/(d^m)\), of dimension \(4m-4\). Surjectivity follows
from the \(e_i\). The kernel statement follows by divisibility by the
four distinct linear factors. Choose any specified unit \(u\in A_h\) and
write \(U=M_u\). Its four constant Taylor coefficients are nonzero;
conversely that condition makes it a unit by the finite inverse of each
truncated Taylor series. Then \(a_h=\operatorname{ev}U\) is onto with
kernel \(J_h\). The old programme choice \(u=j_h(g/h)\), wherever that
entire quotient defines a unit, is retained exactly by this
construction; it is not substituted by 1.

Extend \(A_h,J_h,U\) and \(a_h\) coefficientwise as in FSR23. If
\(u_{i,b}\) denotes its Taylor coefficient, the full operator remains \[
U_L\bigl((x_{a,i,b})_{a,i,b}\bigr)_{a,i,b}
=\sum_{c=0}^{b}u_{i,c}x_{a,i,b-c}.
\tag{FSR26}
\] Thus no amplitude derivative, multiplicity, or support sector is
omitted. Multiplication by \(s\) still acts on each basis element by
\(\alpha_i\) times that element plus the next jet, with the last next
jet zero. The kernel of the extended \(a_h\) is \(C_L\otimes J_h\),
while the joint kernel of all branch maps on the entire algebra is
\(K_L\otimes A_h\). These are distinct explicitly defined subspaces.

\subsection{6. The complete support-valued finite Weil
pairing}\label{the-complete-support-valued-finite-weil-pairing}

Let \(\sigma\) exchange 1 with 2 and 3 with 4. On \(\mathbb C^4\) set \[
b(c,c')=m\sum_{i=1}^4\overline{c_{\sigma i}}c'_i,
\qquad B_h(x,y)=b(a_hx,a_hy).
\tag{FSR27}
\] Complex conjugation and reindexing by \(\sigma\) prove Hermitian
symmetry. The matrix of \(b\) is the sum of two blocks
\(m\left(\begin{smallmatrix}0&1\\1&0\end{smallmatrix}\right)\), which is
nonsingular and has one positive and one negative direction per block.
Since \(a_h\) is onto with kernel \(J_h\), the radical of \(B_h\) is
\(J_h\) and its inertia is \((2,2,4m-4)\). A full diagonal basis is \[
\frac{U^{-1}(e_1+e_2)}{\sqrt{2m}},\quad
\frac{U^{-1}(e_1-e_2)}{\sqrt{2m}},\quad
\frac{U^{-1}(e_3+e_4)}{\sqrt{2m}},\quad
\frac{U^{-1}(e_3-e_4)}{\sqrt{2m}},
\] together with \(U^{-1}e_i(s-\alpha_i)^b\) for \(1\le i\le4\) and
\(1\le b<m\). Direct evaluation gives squared values \(+1,-1\) and 0 as
stated. This proves the finite result independently of any analytic zeta
assertion.

Give \(C_L\) the involution conjugating coefficients and fixing each
\(E_a\). For \(x=\sum_a E_a\otimes x_a\) and \(y=\sum_a E_a\otimes y_a\)
define \[
\mathbb B_L(x,y)=\sum_{a\neq0}E_a B_h(x_a,y_a).
\tag{FSR28}
\] The definition is unambiguous by FSR14. It is \(C_L\)-Hermitian and
sesquilinear by the scalar identities and orthogonality of the \(E_a\).
Its radical is exactly \(C_L\otimes J_h\): if its values against all
\(y\) vanish, choosing \(y\) supported on one sector shows that the
corresponding \(x_a\) lies in the scalar radical. For every branch
\(j\), \[
\eta_j\bigl(\mathbb B_L(x,y)\bigr)
=B_h((\eta_j\otimes\mathrm{id})x,(\eta_j\otimes\mathrm{id})y).
\tag{FSR29}
\] This is the proved comparison with the earlier scalar packet form.

The intrinsic positive cone of the finite product algebra \(C_L\) is \[
C_L^+=\left\{\sum_a c_a E_a:c_a\in\mathbb R_{\ge0}\right\}.
\tag{FSR30}
\] It equals \(\{v^*v:v\in C_L\}\), since each coordinate is an absolute
square, and conversely each nonnegative coordinate has its nonnegative
real square root. The regular-representation trace is
\(\operatorname{tr}_{\rm reg}(\sum_a c_a E_a)=\sum_a c_a\):
multiplication by \(E_a\) is the projection onto its one-dimensional
coordinate. Consequently \(\operatorname{tr}_{\rm reg}\mathbb B_L\) has
inertia \[
\bigl(2(|L|-1),\ 2(|L|-1),\ (4m-4)(|L|-1)\bigr).
\tag{FSR31}
\] This follows from the direct sum of the explicit diagonal bases after
FSR27. No claim identifies this finite trace with the adelic trace
without the corresponding analytic map.

Retaining only the sum of Boolean-branch traces means using
\(\sum_{j\in J(L)}\eta_j\). The radical of that form is \[
\left(\bigoplus_{j\in J(L)}E_j\otimes J_h\right)
\oplus\left(K_L\otimes A_h\right),
\tag{FSR32}
\] and its inertia is \[
\left(2|J(L)|,\ 2|J(L)|,\ (4m-4)|J(L)|+4m(|L|-1-|J(L)|)\right).
\tag{FSR33}
\] Each observed sector has exactly the scalar radical and each
unobserved sector has its entire \(4m\)-dimensional space in the
radical; this proves both formulas.

For \(B^2\) let \(w\) be FSR16 and choose the fully specified element \[
x=w\otimes U^{-1}(e_1-e_2).
\tag{FSR34}
\] Both branch images of \(x\) vanish, but \[
\mathbb B_{B^2}(x,x)=-2m\,w\neq0.
\tag{FSR35}
\] The equality follows from FSR27 and \(w^2=w\). Thus zero results from
all separate branch observations do not make the full support-valued
form zero or positive. Replacing the minus by a plus in FSR34 gives
\(+2m w\), so the omitted space carries both signs in this finite packet
model. This is a concrete comparison result about the defined form. It
is not a counterexample to RH: the arbitrary \(\rho\) in FSR25 has not
been asserted to be a zeta zero.

Finally the scalar inclusion \(x\mapsto u_1\otimes x\) obeys \[
\mathbb B_L(u_1\otimes x,u_1\otimes y)=u_1 B_h(x,y).
\tag{FSR36}
\] It therefore preserves the old calculation as a diagonal subspace.
Full support adds exact sectors and additional observations; it does not
by itself repair a negative scalar value. This specifies the part of the
old Weil construction that survives and the additional calculation
required for an actual full-support adelic realization.

\end{document}
